\chapter{Introduction}

To the hippocampus is attributed a fundamental role in the encoding, consolidation, and recall of declarative
memories~\cite{eichenbaumHippocampus1999, liaoLearning2024}. To this end, the hippocampus receives afferent inputs from multiple
cortical association areas, where information from diverse sensory modalities is processed. This information converges through the
parahippocampal region, which comprises the perirhinal, postrhinal, and, most importantly, the entorhinal cortex
(EC)~\cite{eichenbaumCortical2000}. This rich convergence of information positions the hippocampus ideally to encode associations
between events, consolidating its central role in learning and memory~\cite{henkeModel2010, berdugo-vegaSharpening2023}.

In a simplified and traditional account, the structure of the hippocampus has been described as consisting primarily of the
trisynaptic circuit: information flows in a single direction from the EC to the dentate gyrus (DG), to CA3, and to CA1. The EC
serves as the main input to the DG via the perforant path (PP), whose neurons connect primarily with Granule Cells (GCs), the
principal excitatory cells of the DG, in a ratio of $1:10$; it is worth noting that the EC also sends afferent inputs to CA3 and
CA1, which challenges the simplified notion of the trisynaptic circuit~\cite{basuCorticohippocampal2015}. Precisely because of
this divergent EC-to-DG connectivity, it is believed that the DG is responsible for a computational function called pattern
separation, through its sparse activity and the spike rate of GCs~\cite{hainmuellerDentate2020, kesnerMnemonic2006,
yassaPattern2011}. Pattern separation disambiguates similar information, enabling the identification and differentiation of
distinct yet highly similar experiences, which is fundamental for avoiding catastrophic interference between memories stored in
CA3~\cite{rollsMechanisms2013}.

To illustrate this computation, consider the following oversimplified example: experiences from the same place, with the same
people and objects present and following a similar routine, are likely encoded in the EC in a very similar manner; however, owing
to pattern separation performed by the DG, small details such as a specific conversation or a news item on television may be
sufficient to elicit a neural activity pattern in the DG that is very distinct from those of past experiences, allowing the
encoding of a new memory~\cite{eichenbaumHippocampus2004}.

In contrast to the divergent projection from the EC to the DG, the projection from the DG to CA3 is highly convergent. The CA3
region is characterized by a network of recurrent connections, where its pyramidal CA3 cells (pCA3) excite both interneurons and
other pCA3 cells. This architecture forms an autoassociative network that allows the storage of patterns as neural assemblies:
groups of neurons that represent the same pattern by spiking together, formed through Spike-Timing-Dependent Plasticity (STDP),
which potentiates the connections of neurons that spike together~\cite{kopsickFormation2024}. Once an assembly is formed through
the strengthening of synapses between its cells, it can subsequently be reactivated by a partial input signal from the original
pattern, since activating a subpopulation of assembly cells will cause the remaining cells to become active due to the
associations formed by synaptic potentiation. This process is known as pattern completion, and recent computational models
demonstrate that this architecture allows robust memory retrieval even from degraded cues~\cite{kopsickFormation2024,
leduigouRecurrent2014}. Via the Schaffer collaterals, CA3 projects to CA1, which acts as an interface between the hippocampus
and the neocortex, where memories are consolidated and recalled~\cite{bartschCA12011}.

Beyond its fundamental role in pattern separation, the DG is one of the only areas of the mammalian brain to exhibit adult
neurogenesis, the formation of new neurons from stem cells following development, continuously generating Granule Cells throughout
life~\cite{boldriniHuman2018, dumitruIdentification2025}. Adult neurogenesis in mammals occurs only in a few other specific
areas, such as the olfactory bulb in rodents and possibly the striatum, amygdala, and hypothalamus in
humans~\cite{alonsoImpact2024, jurkowskiHippocampus2020}.

The GCs generated by adult neurogenesis take several weeks to fully develop in rodents. At 3 weeks, the vast majority of
immature Granule Cells (iGCs) die, and from approximately 8 weeks onward they become practically indistinguishable from mature
Granule Cells (mGCs)~\cite{denoth-lippunerFormation2021}. Over the course of maturation, iGCs undergo various morphological and
electrophysiological changes, expanding their dendrites and forming connections with CA3. A critical period occurs at
approximately 4--6 weeks, during which iGCs, already partially integrated into the circuit, exhibit greater excitability and
synaptic plasticity compared to mGCs~\cite{zhaoDistinct2006, denoth-lippunerFormation2021, aimoneRegulation2014}, offering a
critical window for plasticity and memory encoding~\cite{berdugo-vegaSharpening2023}.

In general, adult neurogenesis is implicated in all phases of memory, including encoding, consolidation, recall, and even
forgetting~\cite{chavanMemory2025}. iGCs play a crucial role in novelty detection and pattern separation, enabling the
discrimination between similar memories. Furthermore, recent evidence suggests a specialized role for these neurons in memory
consolidation during REM sleep~\cite{chavanMemory2025} and in the consolidation process that generalizes memories, being
essential for structural changes in CA3 circuits through the non-selective activation of neural assemblies via their high
excitability~\cite{koSystems2025}.

Despite numerous studies demonstrating the role of adult neurogenesis in different hippocampal
functions~\cite{chavanMemory2025, berdugo-vegaSharpening2023}, the computational mechanisms at the circuit level through which
iGCs perform these functions remain poorly understood.

This central question unfolds into two main theoretical lines regarding the function of
iGCs~\cite{berdugo-vegaSharpening2023}. The first hypothesis is that iGCs participate directly in memory encoding, integrating
directly into long-term memory engrams and fulfilling a role similar to mGCs, but with their own peculiarities. In this
scenario, they would act in the encoding of new information, in pattern integration, or in the prevention of catastrophic
interference with memories already established by mGCs. The second hypothesis, which has been explored more recently, proposes
that iGCs function as temporary circuit-modulating elements. In this role, they could both regulate DG activity, increasing the
sparsity of activation by inhibiting mGCs, and modulate CA3 activity~\cite{aimoneComputational2016,
berdugo-vegaSharpening2023}, rather than being directly recruited in memory encoding as mGCs are.

Given the difficulty of experimentation in the DG and the complexity of information processed in the hippocampus, computational
models have been invaluable for understanding the role of neurogenesis~\cite{aimoneComputational2016}. Computational models vary
in scale, biological plausibility, the circuits modeled, and various other factors; therefore, a range of different results
exists regarding the function and mechanisms of neurogenesis~\cite{aimoneComputational2016, berdugo-vegaSharpening2023}.

In a conductance-based DG model~\cite{kimEffect2024}, the large population of mGCs performed excellent pattern separation, while
the minority of iGCs served as pattern integrators. The presence of iGCs deteriorated the DG's pattern separation capacity
compared to a DG composed solely of mGCs; however, the authors postulate that this mixed coding of pattern integration and
separation could increase the capacity for memory storage and retrieval, a hypothesis they were unable to test since only the DG
was modeled, without CA1 or CA3. Interestingly, these results contradict experimental studies demonstrating that an increase in
neurogenesis enhances pattern separation~\cite{sahayIncreasing2011}, which may indicate that computational models of the
hippocampus have not yet fully captured the computations it performs.

In contrast, the work of~\citeonline{yangDynamic2025} employed a DG model to investigate the dynamic impact of neurogenesis on
pattern separation. Their results indicate that the general effect of iGCs on GCs is inhibitory, which increases network
sparsity and consequently improves pattern separation, providing computational evidence for the ``indirect coding'' hypothesis.
The authors also investigated how iGC characteristics during maturation, such as decreasing activity and increasing synaptic
plasticity, dynamically affect pattern separation. The study further suggests that the presence of iGCs increases the
adaptability of the DG network to different input frequencies, improving the robustness of pattern separation.

The model of~\citeonline{aimoneComputational2009} was implemented at an intermediate level of biological
plausibility\footnote{This model is considered to be at an intermediate level of biological plausibility because it does not
model the electrical dynamics of individual neurons through conductance-based models, as in this project, but instead uses a
digitized spike-rate model with a time step of \SI{25}{\milli\second}.}, which allowed for the investigation of iGCs throughout
their entire maturation over time. By modeling the maturation timeline of this iGC population, the model showed that iGCs
preferentially encoded information received when they were young, thus exhibiting the role of temporal integrators. This model
was corroborated by experimental studies~\cite{berdugo-vegaSharpening2023}, but did not analyze the implications for other
areas such as CA1 and CA3, modeling only the DG.

This master's project proposes the development of a computational model of the DG-CA3 circuit to investigate the function of
adult neurogenesis. Unlike many works that focus solely on the DG or on static populations of immature cells, this study will
adopt a more holistic and thus innovative approach. The model will simulate the temporal maturation of iGCs into mGCs, enabling
a dynamic analysis of how the continuous integration of new neurons affects the circuit over time. The main objective is to
analyze the impacts of neurogenesis not only on DG pattern separation, but also on the autoassociation and pattern completion
functions of the CA3 region. The effects of different levels of iGC integration, their heightened excitability, and how
CA3-to-DG feedback projections, a feedback mechanism frequently neglected in models~\cite{myersPattern2011}, influence network
dynamics will be investigated. In doing so, the aim is to understand how iGCs, during their immature phase and after maturation,
modulate the formation and reactivation of memories in CA3.
